\chapter{Estado del arte}\label{cap:estado-del-arte}

Este capítulo sitúa SavePoint en el estado del arte que sostiene el trabajo. Examina primero
las fuentes de datos de videojuegos disponibles y sus condiciones de uso; después, los sistemas
de recomendación; y termina con las metodologías de desarrollo asistido por agentes de código.
Los productos de referencia con los que se compara SavePoint se presentan en la
Sección~\ref{sec:contexto-problema}.

\section{Fuentes de datos de videojuegos}\label{sec:arte-datos}

Tres fuentes públicas cubren la mayor parte del catálogo de videojuegos: Wikidata, IGDB y
\textit{RAWG Video Games Database} (RAWG). Difieren sobre todo en su licencia y en sus condiciones de reutilización, que son
determinantes para un trabajo que debe citar fuentes defendibles. Wikidata publica sus datos
estructurados bajo la dedicación al dominio público \textit{Creative Commons Zero} (CC0), la
licencia más limpia de las tres, aunque el modelado de género y plataforma para videojuegos es
desigual y depende de la edición voluntaria~\cite{wikidata-licensing}. Sus archivos
multimedia viven en Wikimedia Commons y no comparten esa licencia: cada imagen lleva su propia
condición, que puede exigir crédito, de modo que su reutilización requiere una revisión por
archivo~\cite{commons-reuse}. IGDB, propiedad de Twitch, ofrece una interfaz de programación de aplicaciones (API) de versión 4 con una
entidad de plataforma de primera clase y se sirve bajo el \textit{Twitch Developer Services
Agreement} (DSA)~\cite{igdb-api, twitch-dsa}. RAWG ofrece una cobertura comparable, con
metadatos de tienda, pero su nivel gratuito es más restrictivo en cuota y en atribución~\cite{rawg-api}.
La Tabla~\ref{tab:data04}, en la Sección~\ref{cap:datos}, compara las tres fuentes y recoge la
elección de Wikidata para el corpus curado inicial y de IGDB para el catálogo a escala real.

\section{Sistemas de recomendación}\label{sec:arte-recomendacion}

Los sistemas de recomendación se agrupan de forma clásica en tres familias~\cite{ricci2022handbook,
aggarwal2016recommender}. Los métodos basados en contenido describen cada ítem con un vector
de rasgos, como género, plataforma o etiquetas, para recomendar ítems parecidos a los que la
persona ha valorado bien; su ventaja es que funcionan con un solo usuario y explican bien sus
resultados, mientras que su límite es la endogamia, es decir, la tendencia a recomendar
siempre más de lo mismo. Los métodos de filtrado colaborativo ignoran el contenido y explotan
la matriz de interacciones entre personas e ítems, con la hipótesis de que quienes
coincidieron en el pasado volverán a coincidir~\cite{adomavicius2005toward}; capturan gustos
que el contenido no expresa, pero sufren el problema del arranque en frío cuando un usuario o
un ítem no tienen historial. Los métodos híbridos combinan ambas señales para mitigar las
debilidades de cada una~\cite{burke2002hybrid}. SavePoint implementa y compara, sobre su propio
catálogo, representantes de las tres familias, con la formulación y la comparación con datos
reales que desarrollan la Sección~\ref{cap:algoritmos} y el Capítulo~\ref{cap:resultados}.

\section{El método \textit{Get Stuff Done}: agentes coordinados sobre un contexto compartido}\label{sec:arte-agentes}

\textit{Get Stuff Done} (GSD) es un método de trabajo que organiza el desarrollo en agentes
especializados coordinados a través de un contexto compartido y versionado, en lugar de a
través de una conversación continua entre una persona y una única herramienta. La observación
de partida es sencilla: una tarea larga tiende a perder coherencia entre lo decidido al
principio y lo decidido al final, sobre todo cuando se turnan varios agentes para avanzarla y
esa pérdida de coherencia es tanto más probable cuanto más depende de que alguien recuerde el
historial completo de lo hecho hasta el momento~\cite{anthropic-context}. GSD responde a esa
observación con un principio de diseño simple: aislar cada incremento de trabajo en una tarea
atómica, ejecutar cada tarea con el contexto limpio y hacer que la información necesaria para
continuar viaje entre tareas a través de ficheros versionados en el propio repositorio, nunca a
través del historial de una conversación.

Sobre ese principio, GSD organiza el desarrollo en incrementos que recorren un ciclo ordenado~\cite{gsd}. Lo interesante del método no
es tanto ese ciclo, habitual en cualquier proceso de ingeniería disciplinado, como la forma en
que lo instrumenta: cada incremento se apoya en un árbol de ficheros compartido que agentes con
roles distintos leen y escriben de forma coordinada, cada uno acotado a su función
y sin heredar el razonamiento de los demás. El
Capítulo~\ref{cap:gestion} describe con detalle cómo se aplicó este enfoque a SavePoint, qué
mecanismos concretos lo instrumentan y qué evidencia dejó.
